\documentclass[a4paper,12pt,oneside,openany]{report}  
\input{FormatoPack}

\begin{document}

\begin{center}
\textbf{UNIVERSIDADE FEDERAL DO RIO DE JANEIRO}
\vspace{-0.2cm}

\textbf{ESCOLA POLITÉCNICA}
\vspace{-0.2cm}

\textbf{DEPARTAMENTO DE ENGENHARIA ELETRÔNICA E DE COMPUTAÇÃO}
\vspace{0.8cm}

\underline{\textbf{PROPOSTA DE PROJETO DE GRADUAÇÃO}}

Aluno: Lucas Garcia Santiago de Abreu
\vspace{-0.2cm}

lucasgarcia@poli.ufrj.br 

Orientadores: Pedro Henrique González Silva e Amanda Ferreira de Azevedo. 

\vspace{0.2cm}

Curso: ( ) Eletrônica e Computação;  (X) Computação e Informação.
\end{center}


\textbf{1. TÍTULO}

Otimização Estocástica Aplicada ao Gerenciamento de Portfólios de Energia: Uma Abordagem via Programação Dinâmica Dual Estocástica (SDDP)

\vspace{0.4cm}
\textbf{2. ÊNFASE}

Computação e Informação.

\vspace{0.4cm}
\textbf{3. TEMA}

O tema do trabalho é a modelagem e resolução do problema de comercialização de energia elétrica sob incerteza. O foco é otimizar o gerenciamento do portfólio financeiro de uma empresa comercializadora com portfólio próprio, utilizando algoritmos de decomposição avançados para lidar com a alta dimensionalidade estocástica do mercado de energia.

\vspace{0.4cm}
\textbf{4. DELIMITAÇÃO}

O objeto de estudo é a formulação de um modelo de programação linear multiestágio para a decisão de \textit{trading} e controle de fluxo de caixa corporativo. O trabalho está delimitado à modelagem estocástica do Preço de Liquidação de Diferenças (PLD) e da geração física de forma neutra ao risco. A abordagem tradicional via Equivalente Determinístico (DEQ) será utilizada estritamente como linha de base para validação matemática da regra de negócio e para justificar, de forma empírica, a necessidade algorítmica do SDDP frente aos limites de hardware. 

\vspace{0.4cm}
\textbf{5. JUSTIFICATIVA}

O planejamento tático no setor elétrico lida com variáveis altamente estocásticas. Para que uma empresa comercializadora maximize seus resultados e garanta a integridade do seu limite de crédito ao longo do tempo, é imperativo o uso de modelos de otimização sob incerteza. Contudo, a modelagem tradicional (DEQ) é severamente limitada por conta da dimensionalidade elevada que caracteriza o problema, tornando-se intratável em janelas de planejamento realistas devido ao consumo exaustivo de memória RAM. A adoção da Programação Dinâmica Dual Estocástica (SDDP) justifica-se por contornar esse gargalo, dividindo o problema em subproblemas eficientes. O trabalho tem forte apelo prático ao aplicar um algoritmo de ponta para solucionar um desafio latente e complexo do mercado brasileiro.

\vspace{0.4cm}
\textbf{6. OBJETIVO}

O objetivo geral é estudar e implementar um modelo de otimização estocástica multiestágio aplicado ao gerenciamento de portfólios de comercialização de energia elétrica.

Desta forma, os objetivos específicos são: 
\begin{enumerate}[(1)]
    \item Formular a dinâmica físico-financeira do problema em um modelo matemático;
    \item Implementar o modelo na linguagem Julia via DEQ para estabelecer uma prova de conceito e demonstrar sua limitação de escalabilidade (teste de estresse computacional);
    \item Implementar e solucionar o problema completo utilizando o pacote \texttt{SDDP.jl} para superar os gargalos computacionais e obter a política ótima de longo prazo.
\end{enumerate}

\vspace{0.4cm}
\textbf{7. METODOLOGIA}

O trabalho utilizará dados públicos do histórico da Câmara de Comercialização de Energia Elétrica (CCEE) para calibrar um processo estocástico autorregressivo (PAR(1)), responsável por gerar os cenários sintéticos de PLD e geração.

A modelagem do sistema físico e financeiro da empresa será escrita na linguagem Julia. O modelo restrito (DEQ) será implementado através da biblioteca JuMP, e o modelo em larga escala utilizará a biblioteca \texttt{SDDP.jl}. A otimização linear de ambos usará o \textit{solver} HiGHS.

A validação da metodologia incluirá a comparação formal entre os dois modelos para instâncias pequenas (garantindo a equivalência matemática), seguida de um teste de estresse de hardware. Após validar a intratabilidade do DEQ, o algoritmo SDDP será utilizado para simular a política ótima da mesa de operações para o horizonte completo, analisando as decisões de compra e venda nos cenários estocásticos.

\vspace{0.4cm}
\textbf{8. MATERIAIS}
      
Computador pessoal, ambiente de programação Julia (pacotes JuMP, SDDP.jl, DataFrames), ambiente Python (Pandas, Numpy), \textit{solver} HiGHS, e base de dados histórica de PLD e geração da CCEE.

\vspace{0.4cm}
\textbf{9. CRONOGRAMA}

Apresentada graficamente na Figura \ref{Fig:Cronograma}. O trabalho teve início em Março de 2026, com previsão de conclusão em Julho do mesmo ano.

Fase 1: Levantamento bibliográfico e coleta de dados (CCEE).

Fase 2: Implementação do gerador de cenários sintéticos (Python).

Fase 3: Modelagem e implementação da regra de negócio em DEQ.

Fase 4: Comparação matemática, teste de estresse computacional e validação do limite do DEQ.

Fase 5: Implementação do algoritmo SDDP e extração dos resultados do portfólio completo.

Fase 6: Escrita e formatação final do Trabalho de Conclusão de Curso.

\begin{figure}[h]
\centering
\begin{ganttchart}[
    x unit=1.2cm,
    y unit title=0.5cm,
    y unit chart=0.7cm,
    vgrid,
    hgrid,
    title height=1,
    title label font=\bfseries\footnotesize,
    bar/.style={fill=gray!80},
    bar height=0.7,
    progress label text={},
    title/.style={fill=none}
]{1}{5} 
  \gantttitle{2026}{5} \\
  \gantttitle{Mar}{1} \gantttitle{Abr}{1} \gantttitle{Mai}{1} \gantttitle{Jun}{1} \gantttitle{Jul}{1} \\
  
  % Fase 1: Levantamento bibliografico (Marco a Abril)
  \ganttbar{Fase 1}{1}{2} \\
  
  % Fase 2: Geracao de Cenarios Python (Marco a Abril)
  \ganttbar{Fase 2}{1}{2} \\
  
  % Fase 3: Modelo DEQ e Equivalencia (Abril a Maio)
  \ganttbar{Fase 3}{2}{3} \\
  
  % Fase 4: Teste de Hardware e Colapso (Maio a Junho)
  \ganttbar{Fase 4}{3}{4} \\
  
  % Fase 5: Modelo SDDP e Resultados Taticos (Junho a Julho)
  \ganttbar{Fase 5}{4}{5} \\
  
  % Fase 6: Monografia e Entrega (Julho)
  \ganttbar{Fase 6}{5}{5}
\end{ganttchart}
\caption[\small{Cronograma de atividades (Janeiro a Julho/2026)}]{\label{Fig:Cronograma} \footnotesize{Cronograma de atividades planejado.}}
\end{figure}
\begin{thebibliography}{}

\bibitem{Dowson21} DOWSON, O.; KAPELEVICH, L. "SDDP.jl: a Julia package for stochastic dual dynamic programming". INFORMS Journal on Computing, 2021.

\bibitem{CCEE} CCEE. "Dados Históricos de Preços e Geração". Disponível em: http://www.ccee.org.br. Acesso em: 2026.

\end{thebibliography}

      \vspace{2cm}
      \noindent
Rio de Janeiro, 2 de abril de 2026

      \vspace{0.5cm}
      \begin{flushright}
         \parbox{13cm}{
            \vspace{3cm}
            \hrulefill

            \vspace{-.375cm}
            \centering{Lucas Garcia Santiago de Abreu - Aluno}

            \vspace{3cm}
            \hrulefill

            \vspace{-.375cm}
            \centering{Pedro Henrique González Silva - Orientador}
 
            \vspace{2cm}
         }
      \end{flushright}
      \vfill
      
\end{document}